\subsection*{Axioms}
\label{sec:axiom-systems-axioms}

\begin{toolkitbox}{Axioms Toolkit --- Quick Reference}
\small
\begin{tabular}{l l l}
\toprule
\textbf{Concept} & \textbf{Role} & \textbf{Detail} \\
\midrule
Axiom & formal assertion adopted without proof & \hyperref[def:axiom]{Definition} \\
Axiom system & collection of assertions in one language & \hyperref[def:axiom-system]{Definition} \\
Assumption & temporary assertion for reasoning & \hyperref[def:assumption]{Definition} \\
Consistency & absence of contradiction & forward reference \\
Theory & deductive closure of axioms & forward reference \\
\bottomrule
\end{tabular}
\end{toolkitbox}

\begin{axiombox}{Axiom Presentation: A Simple Arithmetic Axiom System}
Reuse the arithmetic signature
\[
  \Sigma_{\mathrm{arith}}
  =
  \bigl(
    \{0,1\},
    \{+,\cdot,S\},
    \{<\},
    \operatorname{ar}
  \bigr)
\]
from the Signatures section. The signature supplies the symbols. An axiom
system supplies assertions formed with those symbols.

For example, the following are familiar arithmetic-style axioms:
\[
  \forall x\;(x+0=x),
\]
\[
  \forall x\;(S(x)\neq 0),
\]
\[
  \forall x\forall y\;
  (S(x)=S(y)\rightarrow x=y).
\]

These examples are illustrative. The signature says that $0$, $+$, and $S$
are available symbols. The axioms assert formal relationships involving those
symbols.
\end{axiombox}

\begin{remark*}[Structural remarks]
\textbf{Vocabulary vs assertion.}
A signature is vocabulary: it specifies what may be said. An axiom is an
assertion: it specifies what is adopted as true within that vocabulary. The
signature supplies symbols; the axioms supply assertions involving those
symbols.

\textbf{Same signature, different axiom systems.}
A single signature may support multiple collections of axioms. For example,
geometric languages can be used for Euclidean or non-Euclidean axiom systems.
The usual group signature can be used for group theory, while adding the
commutativity axiom gives the stronger axiom system for abelian groups. The
vocabulary can remain fixed while the assertions change.

\textbf{Forward conceptual roadmap.}
The chapter proceeds from collections of adopted assertions to their deductive
and semantic roles:
\[
\begin{array}{c}
\text{Axiom System}\\
\Downarrow\\
\text{Theory}\\
\Downarrow\\
\text{Models}
\end{array}
\]
The present section introduces the assertion layer. Later sections formalize
the theory generated by axioms and the models in which those assertions hold.
\end{remark*}

\begin{definition}[Axiom]
\label{def:axiom}
An \emph{axiom} is a formal statement adopted without proof within a specified
language.
\end{definition}

\begin{remark*}[Interpretation]
An axiom is not part of the vocabulary itself. It is an assertion made using
the vocabulary of a language. The signature determines which symbols may
appear; the axiom states something with those symbols.
\end{remark*}

\begin{dependencies}
\begin{itemize}
  \item \hyperref[def:language]{Language}
\end{itemize}
\end{dependencies}

\begin{definition}[Axiom System]
\label{def:axiom-system}
An \emph{axiom system} is a collection of axioms expressed in a common
language.
\end{definition}

\begin{remark*}[Interpretation]
An axiom system records the assertions chosen for a mathematical setting. Two
axiom systems may use the same language while adopting different assertions.
The shared language fixes the vocabulary; the axiom system fixes the adopted
formal claims.
\end{remark*}

\begin{dependencies}
\begin{itemize}
  \item \hyperref[def:axiom]{Axiom}
  \item \hyperref[def:language]{Language}
\end{itemize}
\end{dependencies}

\begin{definition}[Assumption]
\label{def:assumption}
An \emph{assumption} is a statement temporarily accepted for purposes of
reasoning.
\end{definition}

\begin{remark*}[Interpretation]
An axiom is adopted as part of the formal background of a system. An
assumption is local to an argument: it may be introduced to prove a conditional,
to reason within a case, or to derive a contradiction. Axioms belong to the
system; assumptions belong to the proof context.
\end{remark*}

\begin{dependencies}
\begin{itemize}
  \item \hyperref[def:axiom]{Axiom}
\end{itemize}
\end{dependencies}
